% 396_LZ_Vorabfestlegung_En_ch.tex
% Johann Pascher, 28 Oct 2026
% Advance commitment on the LZ event at 248 keV

\providecommand{\BM}{\textbf{[B]}}
\providecommand{\KM}{\textbf{[K]}}
\providecommand{\SM}{\textbf{[S]}}
\providecommand{\XM}{\textbf{[X]}}
\providecommand{\QM}{\textbf{[Q]}}

\chapter*{Advance Commitment on the LZ Event at 248 keV}
\addcontentsline{toc}{chapter}{Advance Commitment on the LZ Event at 248 keV}

\begin{abstract}
In September 2026 the xenon detector LZ published a single event that, by every criterion checked,
looks like a heavy particle striking a xenon nucleus, with a recoil energy of about 248\,keV, against
an expected background of about one hundredth of an event. The global significance is
2.6\,$\sigma$; according to the collaboration the analysis is not blind. In FFGFT, dark matter is not
a substance but the regime below the acceleration scale $a_0$ in which inertia is nonlinear
(Doc.~308). A direct detection of a dark matter particle is therefore a falsification criterion for
this reading. Before further data are analysed, this document states what FFGFT expects for LZ and
what would refute it: no accumulation of such events and no annual modulation peaking in early
June. A confirmed excess with this modulation would refute the reading. The way out, that such a
particle could exist without carrying the rotation curves, is explicitly not claimed.
\end{abstract}

% ============================================================
\section*{Status markers}
% ============================================================
\begin{center}
\begin{tabular}{@{}lp{11cm}@{}}
\textbf{[B]} & Algebraically proven. \\[2pt]
\textbf{[K]} & Numerically confirmed (check script). \\[2pt]
\textbf{[S]} & Assumption, open, or classification without own calculation. \\[2pt]
\textbf{[X]} & Ruled out. \\[2pt]
\textbf{[Q]} & Taken from the literature, with source. \\
\end{tabular}
\end{center}
Check script: \texttt{2/python/Dok396\_Skripte/pruef\_396\_lz\_vorabfestlegung.py}.

% ============================================================
\section{The event}
% ============================================================

LZ sits almost one and a half kilometres underground in the former Homestake mine in South Dakota
and contains seven tonnes of liquid xenon. A collision in the xenon produces two flashes of light,
one in the liquid and a second one when the released electrons reach the gas layer. The ratio of the
two flashes shows whether an electron or a nucleus was struck; their time difference and the light
distribution give the position of the collision.

For the new analysis the collaboration widened the energy window for nuclear recoils from $55$ to
$270$\,keV and analysed 220 live days with 4.71 tonnes of fiducial mass \QM~\cite{LZ2026}. It
contains a single event from 16 June 2023 with $248\pm23\,(\text{stat})\pm23\,(\text{sys})$\,keV. It
lies in the middle of the fiducial volume, more than 25\,cm from wall and bottom, the outer veto
system recorded nothing, and it lies exactly in the nuclear recoil band, far from the band of
ordinary electron recoils. All known background sources together were expected to give
$0.0106\pm0.0008$ events in this region \QM. The authors examined neutrons, a decay with lost charge,
neutrinos and accidental coincidences; none of these sources explains the event convincingly, above
all because the weaker accompanying collisions that almost every known source would have to leave
are missing \QM.

The collaboration itself names two limitations. First, the analysis is not blind, because the
artificial events mixed in for control did not cover the high-energy region. Second, after
accounting for the 616 model variants tested, the significance is $2.6\,\sigma$ globally and up to
$3.4\,\sigma$ locally for individual models \QM. A chance pattern of this kind appears at
$2.6\,\sigma$ about once in 215 trials; particle physics requires $5\,\sigma$ for a discovery, about
once in 3.5 million \KM. From the count alone, without the position of the event in the band, an
expectation of $0.0106$ events gives a probability of $1.05\,\%$ for at least one \KM.

Among the dark matter interpretations, inelastic dark matter fits best
\QM~\cite{TuckerSmithWeiner2001}: in the collision the particle passes into a slightly heavier state
and pays for this energy out of its motion. Slow particles cannot do this, so the weak collisions are
missing and only a few strong ones remain. For an elastic collision of this strength the particle
would have to weigh at least about 75\,GeV \KM.

% ============================================================
\section{The position of FFGFT}
% ============================================================

The findings that led to dark matter are real: flat rotation curves, gravitational lenses, the
Bullet Cluster. FFGFT does not trace them back to an additional substance but to a force law that
follows from $\tilde T\cdot m=1$ together with the K3 geometry,
\[
  a=\sqrt{a_N^2+a_N\,a_0},
\]
with the acceleration scale $a_0$ (Doc.~308, 365). In this reading, dark matter does not denote a
missing mass fraction but the regime below $a_0$ in which inertia is nonlinear. This is established
for the gas-dominated dwarf galaxy DDO~154 \KM; for spiral galaxies deviations of about 19\,\% remain
open, and for the Bullet Cluster it is not established (Doc.~308, 384) \SM.

Without any further assumption this reading implies: there is no dark matter particle, and therefore
no halo of such particles through which the Earth moves. The book ``Dark Matter -- an Illusion''
therefore names the direct detection of such a particle explicitly as a criterion on which the
reading would fail. The LZ event is exactly a measurement of this kind.

% ============================================================
\section{Advance commitment}
% ============================================================

The following statements apply from 28 October 2026, before the further LZ data have been analysed.

\textbf{First expectation: no accumulation.} In the data not yet analysed, several times the 220
live days, collected at least until 2028, the number of nuclear recoils above about $55$\,keV remains
consistent with the expected background. A second or third event at a background of about $0.01$
per 220 days would already be a clear indication against this expectation \SM.

\textbf{Second expectation: no annual modulation.} If the Earth moves through a halo of dark matter,
the motions of the Earth and the Sun add up in summer. The velocity relative to the halo then varies
by about $\pm6\,\%$ with a maximum around 2 June \KM; for inelastic dark matter, which detects only
the fastest particles, this becomes a strong variation of the rate. The event so far occurred
14 days after this maximum, at 97\,\% of the full phase \KM, which on its own means nothing. FFGFT has
no halo and therefore expects no such modulation \SM.

\textbf{Third expectation: no independent confirmation.} Other xenon detectors that examine the same
energy range find no corresponding excess \SM.

\textbf{Refutation criterion.} A confirmed accumulation of such events with an annual maximum in early
June, or an independent confirmation in a second detector at $5\,\sigma$ or more, refutes the
reading of dark matter as a regime without substance. It would then have to be marked \XM.

\textbf{Scope of the refutation.} Only this one reading would be affected: the force law with $a_0$
as the sole explanation of rotation curves and lensing. The other parts of FFGFT do not depend on it.
The derivations from $\xi$ in the particle sector (lepton masses, $\alpha$, Koide form), gravitation
through $\tilde T\cdot m=1$ in the solar system and the achromatic redshift without expansion remain
untouched; the redshift reading is tested against its own criteria (Doc.~308). A dark matter
particle would then be an additional mode in FFGFT, whose origin in the structure of the carrier
would still have to be shown \SM.

\textbf{No way out afterwards.} A single event fixes only the product of local density and cross
section. One could therefore object that such a particle might exist without carrying the rotation
curves. FFGFT does not claim this objection: it would turn a testable reading into an untestable
one. Whoever admits a new substance as soon as it is measured predicts nothing about that substance.

% ============================================================
\section{What the event does not mean}
% ============================================================

A single event at $2.6\,\sigma$ neither refutes FFGFT nor confirms it. There are interpretations
without dark matter, for example a rare neutron that bypassed the shielding and struck the nucleus
without accompanying scattering, or a process within the xenon nucleus itself \QM. FFGFT says
nothing of its own about the nuclear physics of xenon; each of these interpretations is compatible
with it. The value of this commitment lies in the fact that FFGFT has an expectation here that
differs from standard cosmology with a particle halo, and that will be decided in the coming years
on data that nobody has analysed yet.

% ============================================================
\section{Overview}
% ============================================================
\begin{center}
\small
\begin{tabular}{@{}p{7.6cm}p{1.2cm}p{5.0cm}@{}}
\toprule
\textbf{Statement} & \textbf{Status} & \textbf{Reason} \\
\midrule
Event at 248\,keV, background $0.0106$, $2.6\,\sigma$ global, not blind & \QM & LZ 2026 \\
$2.6\,\sigma$: about 1 in 215; $5\,\sigma$: about 1 in 3.5 million & \KM & check script \\
elastic requires $m\gtrsim75$\,GeV & \KM & kinematics \\
dark matter as regime below $a_0$; DDO~154 & \KM & Doc.~308 \\
spiral galaxies, Bullet Cluster & \SM & Doc.~308, 384 \\
no accumulation, no modulation, no confirmation & \SM & advance commitment \\
confirmed accumulation with June maximum refutes the reading & -- & refutation criterion \\
\bottomrule
\end{tabular}
\end{center}

\begin{thebibliography}{9}
\bibitem{LZ2026} LZ Collaboration, search for nuclear recoils at higher energies within the effective field theory of dark matter, preprint, presented at TeVPA 2026 (1 September 2026).
\bibitem{TuckerSmithWeiner2001} D.~Tucker-Smith, N.~Weiner, \textit{Inelastic dark matter}, Physical Review D \textbf{64} (2001) 043502.
\end{thebibliography}
